% Generated draft: requires mathematical and source review.
\documentclass[11pt]{article}
\usepackage[margin=25mm]{geometry}
\usepackage{fontspec}
\setmainfont{Latin Modern Roman}
\usepackage{amsmath,amssymb,mathtools,unicode-math}
\setmathfont{Latin Modern Math}
\usepackage{enumitem,needspace,xurl,hyperref}
\setlength{\parindent}{0pt}
\setlength{\parskip}{5pt}
\setlength{\emergencystretch}{3em}
\newcommand{\N}{\mathbb{N}}
\newcommand{\Z}{\mathbb{Z}}
\newcommand{\Q}{\mathbb{Q}}
\newcommand{\R}{\mathbb{R}}
\newcommand{\C}{\mathbb{C}}
\newcommand{\F}{\mathbb{F}}
\newcommand{\A}{\mathbb{A}}
\newcommand{\PP}{\mathbb{P}}
\newcommand{\E}{\mathbb{E}}
\newcommand{\eps}{\varepsilon}
\newcommand{\dd}{\,\mathrm{d}}
\begin{document}
\section*{UnsolvedMath ID: 30005275}
\textbf{Source problem code:} catalogue entry
\subsection{Definitions and mathematical statement}
The supplied record does not define \(W_f(a;p)\), specify the coefficient domain of \(f,g\), or quantify what makes a prime large. Consequently, the following explicit formulation is conditional on the usual interpretation by polynomial exponential sums.

For a prime \(p\), let \(\F_p=\Z/p\Z\), and define \(e_p(t)=\exp(2\pi i\widetilde t/p)\), where \(\widetilde t\in\Z\) represents \(t\in\F_p\). For \(h\in\F_p[X]\), define the probability measure and exponential sum
\[
\mu_{h,p}(t)=\frac{1}{p}\#\{x\in\F_p:h(x)=t\},\qquad
S_h(a;p)=\sum_{x\in\F_p}e_p(ah(x)).
\]
With the Fourier convention
\[
\widehat\mu_{h,p}(a)=\sum_{t\in\F_p}\mu_{h,p}(t)e_p(at),
\]
we have \(\widehat\mu_{h,p}(a)=p^{-1}S_h(a;p)\). Thus any normalization \(W_h(a;p)=c_pS_h(a;p)\), with the same nonzero factor \(c_p\) for both polynomials, gives the same equality-of-magnitudes condition. The supplied record does not establish that this is its precise definition of \(W_h\).

\textbf{Question, with two source alternatives.} Determine the algebraic relations between \(f\) and \(g\) forced by equality of their Fourier magnitudes, in either of the following settings:
\begin{enumerate}
\item For one prime \(p\) described in the source as large, and polynomials interpreted over \(\F_p\), assume
\[
|S_f(a;p)|=|S_g(a;p)|\qquad\text{for every }a\in\F_p^{\times}.
\]
Determine what algebraic relations this forces between the polynomials. The required meaning of large, any degree restrictions, and whether the intended conclusion concerns polynomials over \(\F_p\) or lifts to another coefficient domain are unspecified.
\item For fixed polynomials \(f,g\) admitting reductions \(\overline f_p,\overline g_p\) at all sufficiently large primes, assume
\[
\exists P_0\ \forall\text{ primes }p>P_0\ \forall a\in\F_p^{\times},\qquad
|S_{\overline f_p}(a;p)|=|S_{\overline g_p}(a;p)|.
\]
Determine the algebraic relations this forces between the fixed polynomials. Their coefficient domain and the reduction maps must be specified before this becomes an unconditional formulation.
\end{enumerate}
Here the source phrase that \(a\) is coprime to \(p\) is interpreted as \(a\in\F_p^{\times}\). Under the displayed exponential-sum interpretation, the equality at \(a=0\) is automatic because \(S_h(0;p)=p\).

The fragment labeled original statement also contains an equality \(W_f(a;p)=W_g(a;p)\) for all \(a\), in a discussion of equality of probability measures. That retains Fourier phase and is stronger than the magnitude hypothesis in the reconstructed statement. Its surrounding definitions are absent. The fragment contrasts Davenport's question, which remembers only the supports of the measures, with the present question, which discards Fourier phase.

The supplied dated assessment reports a relation theorem under restrictive Sidon and Morse hypotheses, but states neither those hypotheses nor its conclusion. They are therefore not added as assumptions or asserted here as a theorem. The intended class of algebraic relations and any necessary exceptional cases also remain unspecified.
\subsection{Short English statement}
Determine what algebraic relations between two polynomials are forced when their polynomial exponential sums have equal absolute values at every nonzero frequency, either for one large prime or for every sufficiently large prime. The coefficient domain, the definition of \(W_f\), the meaning of large, and the intended algebraic conclusion are missing, so the explicit finite-field formulation is conditional. The supplied partially solved label and restricted Sidon--Morse result require independent mathematical review.
\subsection{Sources}
\begin{itemize}
\item Local UnsolvedMath registry, supplied complete record with permanent ID \(30005275\), title ``Recovering Polynomials from Fourier-Magnitude Data,'' exact source problem code \texttt{catalogue entry}. No individual proposer is supplied; the proposed year is recorded as 2022.
\item Supplied report citation: \emph{Analytic Number Theory (2023)}, Oberwolfach Reports, DOI \texttt{10.4171/owr/2022/50}, pp.\ 2895--2960. \url{https://doi.org/10.4171/owr/2022/50}. The record describes Oberwolfach Reports as documenting workshops and their open-problem discussions.
\item Supplied literature source: Kowalski and Soundararajan, \url{https://doi.org/10.1007/978-3-031-05331-3_13}. The original fragment attributes most reported results to Kowalski and Soundararajan [2022]; no title or precise theorem reference is supplied.
\item Supplied literature source: Oberwolfach Report 50, \url{https://ems.press/content/serial-article-files/46986}.
\item Literature checked at: \texttt{2026-08-22}. This is supplied dated curation, not an independently verified assessment. It labels the problem partially solved and reports a relation theorem under restrictive Sidon and Morse hypotheses while describing the extracted general question as unresolved. The background explicitly requests independent verification of time-sensitive status and cited claims.
\item The registry labels the statement \texttt{corrected\_verified} and says it was reconstructed and checked against the official Oberwolfach report. These are supplied provenance claims; no source was independently visited or verified for this draft.
\end{itemize}
\end{document}
